\StartChapter{Funktionen}

\SubsectionBar{Eigenschaften\ZSFScriptRef{54}}
\ZSFindex{Stetigkeit}

Eine \ZSFkeyword{Funktion} $f: D(f) \to Z(f)$ bildet eine Menge auf einer anderen Menge ab. Sie besitzt einen Definitionsbereich $D(f)$, Zielbereich $Z(f)$ und einen Wertebereich / ein Bild $W(f) = \{f(x) | x \in D(f)\}$.

\begin{tablebox}
\begin{ZSFtable}{C C}
    \ZSFheaderRow \ZSFhead{Gerade Funktion} & \ZSFhead{Ungerade Funktion} \\
    $\displaystyle f(-x) = f(x)$ & $\displaystyle f(-x) = -f(x)$ \\
\end{ZSFtable}
\end{tablebox}

\begin{ZSFlist}
    \ZSFItem{surjektiv}{Jeder Wert im Zielbereich wird angenommen: $Z(f) = W(f)$.}
    \ZSFItem{injektiv}{Jede Horizontale schneidet den Graphen höchstens einmal: $f(x_1) = f(x_2) \Rightarrow x_1 = x_2$.}
    \ZSFItem{bijektiv}{Injektiv und surjektiv. Bijektive Funktionen besitzen eine \ZSFkeyword{Umkehrfunktion / Inverse} mit $W(f^{-1}) = D(f)$ und $D(f^{-1}) = W(f)$.}
\end{ZSFlist}
\ZSFindex{surjektiv}
\ZSFindex{injektiv}
\ZSFindex{bijektiv}
\ZSFindexsee{Inverse}{Umkehrfunktion / Inverse}

\textVorBox{Eine Funktion $f(x)$ ist \ZSFkeyword{stetig} im Punkt $x = \xi$, wenn:}

\begin{formulabox}
    $\displaystyle \ZSFlimAuto{x \to \xi^-} f(x) = f(\xi) = \ZSFlimAuto{x \to \xi^+} f(x)$
\end{formulabox}

\SubsectionBar[sec:trigo_funktionen]{Trigonometrische Funktionen\ZSFScriptRef{58,97,98,99}}
\ZSFindex{Trigonometrische Funktionen}
\ZSFindexsee{Winkelfunktionen}{Trigonometrische Funktionen}

\begin{tablebox}
\begin{ZSFtable}[header=false, font=dense, colsep=tight]{Y{1.3} Z{0.95} Z{0.95} Z{0.95} Z{0.95} Z{0.95} Z{0.95}}
    $\alpha(^\circ)$ & $0^\circ$ & $30^\circ$ & $45^\circ$ & $60^\circ$ & $90^\circ$ & $180^\circ$ \\ 
    $\alpha(\pi)$ & $0$ & $\pi/6$ & $\pi/4$ & $\pi/3$ & $\pi/2$ & $\pi$ \\
    $\sin(\alpha)$ & $0$ & $1/2$ & $\sqrt{2}/2$ & $\sqrt{3}/2$ & $1$ & $0$ \\
    $\cos(\alpha)$ & $1$ & $\sqrt{3}/2$ & $\sqrt{2}/2$ & $1/2$ & $0$ & $-1$ \\
    $\tan(\alpha)$ & $0$ & $\sqrt{3}/3$ & $1$ & $\sqrt{3}$ & $\infty$ & $0$ \\
\end{ZSFtable}
\end{tablebox}

\textVorBox{Der Cosinus ist eine gerade Funktion, der Sinus ist ungerade. Der Tangens ist der Quotient aus Sinus und Cosinus:}

\begin{formulabox}
    $\displaystyle \tan(x) = \frac{\sin(x)}{\cos(x)}$
\end{formulabox}

\textNachBox{\begin{ZSFlist}
    \ZSFItem{$\cos(x)$ ist grösser als $\sin(x)$ auf $(0, \frac{\pi}{4})$ und $(\frac{5\pi}{4}, 2\pi)$}
    \ZSFItem{$\sin(x)$ ist grösser als $\cos(x)$ auf $(\frac{\pi}{4}, \frac{5\pi}{4})$}
    \ZSFItem{$\cos(x)$ ist positiv auf $(0, \frac{\pi}{2})$ und $(\frac{3\pi}{2}, 2\pi)$}
    \ZSFItem{$\sin(x)$ ist positiv auf $(0, \pi)$}
\end{ZSFlist}}

\textVorBox{Der Cosinus und Sinus sind phasenversetzt:}

\begin{formulabox}
    \[\begin{aligned}
        \cos(x \pm \pi/2) &= \mp\sin(x), & \sin(x \pm \pi/2) &= \pm\cos(x) \\
        \cos(x \pm \pi) &= -\cos(x), & \sin(x \pm \pi) &= -\sin(x)
    \end{aligned}\]
\end{formulabox}

\textVorBox{Die Ableitungen lauten:}

\begin{formulabox}
    \[\begin{aligned}
        \frac{d}{dx} \cos(x) &= -\sin(x), & \frac{d}{dx} \arccos(x) &= -\frac{1}{\sqrt{1-x^2}} \\
        \frac{d}{dx} \sin(x) &= \cos(x), & \frac{d}{dx} \arcsin(x) &= \frac{1}{\sqrt{1-x^2}} \\
        \frac{d}{dx} \tan(x) &= \frac{1}{\cos^2(x)}, & \frac{d}{dx} \arctan(x) &= \frac{1}{1+x^2}
    \end{aligned}\]
\end{formulabox}
\textVorBox{Die wichtigsten Identitäten lauten:}

\begin{tablebox}[Grundrelationen \& Euler]
\begin{ZSFtable}[header=false]{C C}
    $\displaystyle \sin^2(x) + \cos^2(x) = 1$ & $\displaystyle \cos^2(x) = \frac{1}{1 + \tan^2(x)}$ \\
    $\displaystyle \cos(x) = \frac{e^{ix} + e^{-ix}}{2}$ & $\displaystyle \sin(x) = \frac{e^{ix} - e^{-ix}}{2i}$ \\ 
\end{ZSFtable}
\end{tablebox}

\begin{tablebox}[Arkus-Identitäten]
\begin{ZSFtable}[header=false]{C C}
    \ZSFspan{C}{2}{$\displaystyle \sin(\arccos(x)) = \cos(\arcsin(x)) = \sqrt{1 - x^2}$} \\
    $\displaystyle \sin(\arctan(x)) = \frac{x}{\sqrt{1+x^2}}$ & $\displaystyle \cos(\arctan(x)) = \frac{1}{\sqrt{1+x^2}}$ \\
\end{ZSFtable}
\end{tablebox}
\ZSFindex{Arkusfunktionen}

\begin{tablebox}[Trigonometrische Umkehrung / Intervalle]
\begin{ZSFtable}[header=false]{L}
    $\sin(x) = y,\; x \in [0, \pi] \Rightarrow x = \arcsin(y) \text{ oder } x = \pi - \arcsin(y)$ \\
    $\cos(x) = y,\; x \in [0, \pi] \Rightarrow x = \arccos(y)$ \\
    $\cos(x) = y,\; x \in [0, 2\pi] \Rightarrow x = \arccos(y) \text{ oder } x = 2\pi - \arccos(y)$ \\
\end{ZSFtable}
\end{tablebox}

\begin{tablebox}[Doppel- \& Halbwinkel]
\begin{ZSFtable}[header=false]{C C}
    $\displaystyle \sin(2x) = 2\sin(x)\cos(x)$ & $\displaystyle \tan(2x) = \frac{2\tan(x)}{1 - \tan^2(x)}$ \\
    \ZSFspan{C}{2}{$\displaystyle \cos(2x) = \cos^2(x) - \sin^2(x) = 1 - 2\sin^2(x)$} \\
    $\displaystyle \sin\left(\frac{x}{2}\right) = \pm\sqrt{\frac{1-\cos(x)}{2}}$ & $\displaystyle \cos\left(\frac{x}{2}\right) = \pm\sqrt{\frac{1+\cos(x)}{2}}$ \\ 
\end{ZSFtable}
\end{tablebox}

\begin{tablebox}[Vielfache \& Additionstheoreme]
\begin{ZSFtable}[header=false]{C C}
    $\displaystyle \cos(3x) = 4\cos^3(x) - 3\cos(x)$ & $\displaystyle \sin(3x) = 3\sin(x) - 4\sin^3(x)$ \\
    \ZSFspan{C}{2}{$\displaystyle \cos(a \pm b) = \cos(a)\cos(b) \mp \sin(a)\sin(b)$} \\
    \ZSFspan{C}{2}{$\displaystyle \sin(a \pm b) = \cos(b)\sin(a) \pm \sin(b)\cos(a)$} \\
    \ZSFspan{C}{2}{$\displaystyle \sin(t)\cos(s) = \frac{1}{2}(\sin(t+s) + \sin(t-s))$} \\
\end{ZSFtable}
\end{tablebox}
\ZSFindex{Additionstheoreme}

\textVorBox{Transformation von trigonometrischen Funktionen:}

\begin{formulabox}
    $\displaystyle a\cos(bx + c) + d$
\end{formulabox}

\textNachBox{\begin{ZSFlist}
    \ZSFItem{Bei $a > 1$ Streckung in $y$-Richtung, bei $0 < a < 1$ Stauchung in $y$-Richtung, Spiegelung an der $x$-Achse bei $a < 0$}
    \ZSFItem{Bei $0 < b < 1$ Streckung in $x$-Richtung, bei $b > 1$ Stauchung in $x$-Richtung}
    \ZSFItem{Verschiebung um $c$ in \ZSFlabel{negative} $x$-Richtung}
    \ZSFItem{Verschiebung um $d$ in \ZSFlabel{positive} $y$-Richtung}
\end{ZSFlist}}

\SubsectionBar{Hyperbolische Funktionen\ZSFScriptRef{60}}

\ZSFkeyword{Hyperbolische Funktionen} sind trigonometrische Funktionen an der Einheitshyperbel.

\begin{tablebox}
\begin{ZSFtable}[header=false]{C C C C}
    $x$ & $\cosh(x)$ & $\sinh(x)$ & $\tanh(x)$ \\
    $0$ & $1$ & $0$ & $0$ \\
\end{ZSFtable}
\end{tablebox}

\textVorBox{Ihre Umkehrfunktionen sind definiert als:}

\begin{formulabox}
    \[\begin{aligned}
        \Arcosh(x) &= \ln\left(x + \sqrt{x^2 - 1}\right) \\
        \Arsinh(x) &= \ln\left(x + \sqrt{x^2 + 1}\right) \\
        \Artanh(x) &= \frac{1}{2}\ln\left(\frac{1+x}{1-x}\right)
    \end{aligned}\]
\end{formulabox}

\textVorBox{Die Ableitungen lauten:}

\begin{formulabox}
    \[\begin{aligned}
        \frac{d}{dx}\cosh(x) &= \sinh(x), & \frac{d}{dx}\sinh(x) &= \cosh(x) \\
        \frac{d}{dx}\Arcosh(x) &= \frac{1}{\sqrt{x^2 - 1}}, & \frac{d}{dx}\Arsinh(x) &= \frac{1}{\sqrt{x^2 + 1}}
    \end{aligned}\]
\end{formulabox}

\textVorBox{Die wichtigsten Identitäten lauten:}

\begin{tablebox}[Euler-Formeln]
\begin{ZSFtable}[header=false]{C C}
    $\displaystyle \cosh(x) = \frac{e^x + e^{-x}}{2}$ & $\displaystyle \sinh(x) = \frac{e^x - e^{-x}}{2}$ \\ 
\end{ZSFtable}
\end{tablebox}

\begin{tablebox}[Grundrelation \& Doppelwinkel]
\begin{ZSFtable}[header=false]{C}
    $\displaystyle \cosh(x)^2 - \sinh(x)^2 = 1$ \\
    $\displaystyle \cosh(2x) = \cosh(x)^2 + \sinh(x)^2$ \\ 
\end{ZSFtable}
\end{tablebox}

\begin{tablebox}[Halbwinkel\ \& Area-Identitäten]
\begin{ZSFtable}[header=false]{C C}
    $\displaystyle \sinh\left(\frac{x}{2}\right) = \pm\sqrt{\frac{\cosh(x)-1}{2}}$ & $\displaystyle \cosh\left(\frac{x}{2}\right) = \sqrt{\frac{\cosh(x)+1}{2}}$ \\
    $\displaystyle \sinh(\Arcosh(x)) = \sqrt{x^2 - 1}$ & $\displaystyle \cosh(\Arsinh(x)) = \sqrt{x^2 + 1}$ \\
\end{ZSFtable}
\end{tablebox}
\ZSFindex{Areafunktionen}

\begin{tablebox}[Additionstheoreme]
\begin{ZSFtable}[header=false]{C}
    $\displaystyle \cosh(a \pm b) = \cosh(a)\cosh(b) \pm \sinh(a)\sinh(b)$ \\
    $\displaystyle \sinh(a \pm b) = \sinh(a)\cosh(b) \pm \cosh(a)\sinh(b)$ \\ 
\end{ZSFtable}
\end{tablebox}
\ZSFindex{Additionstheoreme}

\SubsectionBar{Exponential \& Logarithmusfunktionen\ZSFScriptRef{17,58}}

\ZSFkeyword{Exponentialfunktionen} sind Funktionen der Form $a b^x$. Die Umkehrfunktionen der Exponentialfunktionen sind die \ZSFkeyword{Logarithmusfunktionen} $\log_c(x)$.

\ZSFdanger{Hinweis:} Bei Herrn Steiger entspricht $\log(x) = \log_e(x) = \ln(x)$.

\textVorBox{Transformation für Potenz-, Exponential-, und Logarithmusfunktionen:}

\begin{tablebox}
\begin{ZSFtable}[header=false]{Y{1.3} Z{0.7}}
    \ZSFlabel{Potenzfunktion:} & $a(x + b)^n + c$ \\
    \ZSFlabel{Exponentialfunktion:} & $am^x + c$ \\
    \ZSFlabel{Logarithmusfunktion:} & $a\log(x + b) + c$ \\ 
\end{ZSFtable}
\end{tablebox}

\textNachBox{\begin{ZSFlist}
    \ZSFItem{Bei $a > 1$ Streckung in $y$-Richtung, bei $0 < a < 1$ Stauchung in $y$-Richtung, Spiegelung an $x$-Achse bei $a < 0$}
    \ZSFItem{Verschiebung um $b$ in \ZSFlabel{negative} $x$-Richtung}
    \ZSFItem{Verschiebung um $c$ in \ZSFlabel{positive} $y$-Richtung}
\end{ZSFlist}}

\SubsectionBar{Grenzwerte\ZSFScriptRef{51,52,62,66}}
\ZSFindex{Grenzwertberechnung}

Der \ZSFkeyword{Limes / Grenzwert} entspricht dem Funktionswert, dem sich eine Funktion an einer betrachteten Stelle annähert. Existiert ein Grenzwert, konvergiert die Funktion, andernfalls divergiert sie.

\textVorBox{Die \ZSFkeyword{Grenzwertsätze} lauten:}

\begin{formulabox}
    \[\begin{aligned}
        \ZSFlimAuto{x \to a} (f(x) \pm g(x)) &= \ZSFlimAuto{x \to a} f(x) \pm \ZSFlimAuto{x \to a} g(x) \\
        \ZSFlimAuto{x \to a} (f(x)g(x)) &= \ZSFlimAuto{x \to a} f(x) \cdot \ZSFlimAuto{x \to a} g(x) \\
        \ZSFlimAuto{x \to a} \frac{f(x)}{g(x)} &= \frac{\ZSFlimAuto{x \to a} f(x)}{\ZSFlimAuto{x \to a} g(x)}
    \end{aligned}\]
\end{formulabox}

\begin{ZSFlist}[Vorgehen: Grenzwerte berechnen][ordered]
    \ZSFItem{Konstanten}{Können vor den Grenzwert gezogen werden.}
    \ZSFItem{Substituieren}{Bei gewissen Termen (z.\,B. Argumente von $\sin$/$\cos$) lohnt es sich.}
    \ZSFItem{Bernoulli-L'Hôpital}{Falls $\ZSFlimAuto{x \to \infty} f(x) = 0$ und $\ZSFlimAuto{x \to \infty} g(x) = 0$ (oder $\pm\infty$):}
\end{ZSFlist}
\ZSFindex[LHopital]{L'Hôpital}
\ZSFindexsee[Regel von LHopital]{Regel von L'Hôpital}{L'Hôpital}

\begin{formulabox}
    $\displaystyle \ZSFlimAuto{x \to \infty} \frac{f(x)}{g(x)} = \ZSFlimAuto{x \to \infty} \frac{f'(x)}{g'(x)}$
\end{formulabox}

\textVorBox{Eine \ZSFkeyword{Asymptote} $\sim$ einer Kurve ist eine Funktion, bei der der Abstand zwischen Kurve und Funktion gegen Null strebt:}

\begin{formulabox}
    $\displaystyle \ZSFlimAuto{x \to \infty} (f(x) - g(x)) = 0$
\end{formulabox}

\textNachBox{Ist $f$ eine Asymptote von $g$, so ist auch $g$ eine Asymptote von $f$.}

Die \ZSFkeyword{Landau-Symbole} werden verwendet, um das Wachstum von Funktionen zu beschreiben:\\
\textVorBox{$f$ wächst langsamer als $g$:}

\begin{formulabox}
    $\displaystyle f(x) = o(g(x)) \text{ für } x \to a \Leftrightarrow \ZSFlimAuto{x \to a} \frac{f(x)}{g(x)} = 0$
\end{formulabox}

\textVorBox{$f$ wächst höchstens genauso schnell wie $g$:}

\begin{formulabox}
    $\displaystyle f(x) = \mathcal{O}(g(x)) \text{ für } x \to a \Leftrightarrow \ZSFlimAuto{x \to a} \frac{f(x)}{g(x)} \leq C \in \R$
\end{formulabox}

\textNachBox{$f(x) = o(g(x))$ und $f(x) = \mathcal{O}(g(x))$ können gleichzeitig zutreffen, da $C$ den Wert $0$ annehmen kann.}

\textVorBox{Eine Funktion kann gegen 0 oszillieren:}

\begin{formulabox}
    \[\begin{aligned}
        \ZSFlimAuto{x \to \infty} \left(\frac{\sin(x)}{x}\right) &= 0, & \ZSFlimAuto{x \to \infty} \left(\frac{1}{x\sin(x) + e^x}\right) &= 0
    \end{aligned}\]
\end{formulabox}

\textVorBox{Für $x \to \infty$ gilt:}

\begin{formulabox}
    $\displaystyle e^x > a^x > x^b > \sqrt[c]{x} > \log_d(x) > \text{trig}$
\end{formulabox}

\textVorBox{Die \ZSFkeyword{Eulersche Zahl} $e$ kann mit Grenzwerten dargestellt werden:}

\begin{formulabox}
    \[\begin{aligned}
        e &= \ZSFlimAuto{n \to \infty} \left(1 + \frac{1}{n}\right)^n, & e^x &= \ZSFsumAuto{n=0}{\infty} \frac{x^n}{n!} = \ZSFlimAuto{n \to \infty} \left(1 + \frac{x}{n}\right)^n
    \end{aligned}\]
\end{formulabox}

\begin{tablebox}[Wichtige Grenzwerte]
\begin{ZSFtable}[header=false]{C C}
    $\displaystyle \ZSFlimAuto{x \to \infty} \sqrt[x]{x} = 1$ & $\displaystyle \ZSFlimAuto{x \to 0} \frac{e^x - 1}{x} = 1$ \\
    $\displaystyle \ZSFlimAuto{x \to 0} \frac{a^x - 1}{x} = \ln(a)$ & $\displaystyle \ZSFlimAuto{x \to 1} \frac{\ln(x)}{x - 1} = 1$ \\
    $\displaystyle \ZSFlimAuto{x \to 0} \frac{\ln(a + x) - \ln(a)}{x} = \frac{1}{a}$ & $\displaystyle \ZSFlimAuto{x \to 0} \frac{\sin(x)}{x} = 1$ \\
    $\displaystyle \ZSFlimAuto{x \to 0} \frac{x}{\sin(ax)} = \frac{1}{a}$ & $\displaystyle \ZSFlimAuto{x \to 0} \frac{\sin(ax)}{\sin(x)} = a$ \\
    $\displaystyle \ZSFlimAuto{x \to 0} \frac{\cos(x)-1}{x} = 0$ & $\displaystyle \ZSFlimAuto{x \to 0} \frac{1 - \cos(x)}{x^2} = \frac{1}{2}$ \\
    $\displaystyle \ZSFlimAuto{x \to \infty} x\sin\left(\frac{1}{x}\right) = 1$ & $\displaystyle \ZSFlimAuto{x \to 0} \frac{\tan(x)}{x} = 1$ \\
    $\displaystyle \ZSFlimAuto{x \to \pi/2} \tan(x) = \pm\infty$ & $\displaystyle \ZSFlimAuto{x \to 0} \frac{\arctan(x)}{x} = 1$ \\
    $\displaystyle \ZSFlimAuto{x \to \infty} \arctan(x) = \frac{\pi}{2}$ & $\displaystyle \ZSFlimAuto{x \to 0} \frac{\arcsin(x)}{x} = 1$ \\
\end{ZSFtable}
\end{tablebox}

\SubsectionBar{Inverse}
\ZSFindex{Umkehrfunktion / Inverse}

Die Steigung der Inverse ist die Inverse der Steigung.

\textVorBox{Für die Inverse einer Verkettung gilt:}

\begin{formulabox}
    $\displaystyle \ZSFmhlA{f}(\ZSFmhlB{g}(x)) = y \quad \Rightarrow \quad \left(\ZSFmhlA{f}(\ZSFmhlB{g}(x))\right)^{-1} = \ZSFmhlB{g^{-1}}(\ZSFmhlA{f^{-1}}(y))$
\end{formulabox}
